% =====================================================================
% Paper 50 (standalone): Boundary CFT_3-on-S^3 observables on the
% truncated GeoVac spectral triple.
%
% Three complementary observables on the same wedge KMS state via three
% structurally distinct calculations:
%   (A) spectral-zeta side -> universal F-theorem coefficient
%   (B) wedge KMS state side -> boundary-dimension log scaling
%   (C) modular Hamiltonian side -> CHM integer-spectrum identification
%
% Substantive new structural content:
%   1. Master Mellin engine verified on independent (non-GeoVac-internal)
%      physics observables for the first time
%   2. Scalar/Dirac orthogonal decomposition under M2 + M3 projection
%   3. F-coefficient is in spectral-zeta side (A), NOT in wedge KMS
%      entropy side (B): structural refinement of the framework's
%      modular machinery
%
% Twelfth math.OA standalone in the GeoVac series.
% Siblings: Papers 38, 39, 40, 42, 43, 44, 45, 46, 47, 48, 49.
%
% Status: first draft, May 2026, post-Sprint AdS-Tracks-A+B+C unified
% (2026-05-25). Built from:
%   debug/sprint_ads_track_a_partition_function_audit_memo.md
%   debug/sprint_ads_track_b_entanglement_audit_memo.md
%   debug/sprint_ads_track_c_modular_jlms_audit_memo.md
%   debug/ads_track_a_findings_memo.md
%   debug/ads_unified_synthesis_memo.md
%
% Honest scope: boundary-side CFT_3-on-S^3 reproduction; bulk side
% (JLMS, RT, AdS_4 / H^4 reconstruction) BLOCKED at framework's
% existing infrastructure (Class-1 calibration-external; Sprint RH-B
% 2026-04-17 closed S^3 -> H^3 Wick rotation as clean dead end).
% =====================================================================
\documentclass[11pt,a4paper]{article}

\usepackage[T1]{fontenc}
\usepackage[utf8]{inputenc}
\usepackage{amsmath, amssymb, amsthm, mathrsfs, mathtools}
\usepackage{geometry}
\geometry{margin=1in}
\usepackage{hyperref}
\hypersetup{colorlinks=true, linkcolor=blue!50!black, citecolor=blue!50!black, urlcolor=blue!50!black}
% \usepackage{microtype}  % disabled: MiKTeX font-expansion environmental issue
\usepackage{graphicx}
\usepackage{booktabs}
\usepackage[numbers,sort&compress]{natbib}
\usepackage{xcolor}

\theoremstyle{plain}
\newtheorem{theorem}{Theorem}[section]
\newtheorem{lemma}[theorem]{Lemma}
\newtheorem{proposition}[theorem]{Proposition}
\newtheorem{corollary}[theorem]{Corollary}
\theoremstyle{definition}
\newtheorem{definition}[theorem]{Definition}
\newtheorem{remark}[theorem]{Remark}
\newtheorem{example}[theorem]{Example}

\newcommand{\nmax}{n_{\max}}
\newcommand{\sthree}{S^{3}}
\newcommand{\Hilb}{\mathcal{H}}
\newcommand{\Acal}{\mathcal{A}}
\newcommand{\Tcal}{\mathcal{T}}
\newcommand{\Op}{\mathcal{O}}
\newcommand{\opnorm}[1]{\lVert #1 \rVert_{\mathrm{op}}}
\newcommand{\norm}[1]{\lVert #1 \rVert}
\newcommand{\R}{\mathbb{R}}
\newcommand{\C}{\mathbb{C}}
\newcommand{\N}{\mathbb{N}}
\newcommand{\Z}{\mathbb{Z}}
\newcommand{\Q}{\mathbb{Q}}
\newcommand{\Tr}{\mathrm{Tr}}
\newcommand{\DGV}{D_{\mathrm{GV}}}
\newcommand{\DCH}{D_{\mathrm{CH}}}
\newcommand{\HGV}{\mathcal{H}_{\mathrm{GV}}}
\newcommand{\Kalpha}{K_{\alpha}^{W}}
\newcommand{\rhoW}{\rho_{W}}
\newcommand{\Wedge}{W}
\newcommand{\KPS}{\text{KPS}}
\newcommand{\CHM}{\text{CHM}}

\title{Boundary CFT$_{3}$-on-$\sthree$ observables on the truncated\\
GeoVac spectral triple: partition function, wedge KMS entropy,\\
and Casini--Huerta--Myers modular Hamiltonian}
\author{J.~Loutey\thanks{Independent researcher. \texttt{jloutey@gmail.com}.
This work was developed in an AI-augmented agentic workflow with
Anthropic's Claude. Implementation, internal memos, and bibliographic
collation were performed collaboratively; all theorems, design choices,
and editorial decisions are the author's responsibility.}}
\date{May 2026 (preprint)}

\begin{document}
\maketitle

\begin{abstract}
We exhibit three complementary boundary-side CFT$_{3}$-on-$\sthree$
observables on the truncated GeoVac spectral triple, each computed
from the same wedge KMS state $\rhoW = e^{-\Kalpha}/Z$ via a
structurally distinct operation: \emph{(A)} the universal F-theorem
coefficient $F_{\sthree}$, extracted via the spectral-zeta derivative
$-\tfrac{1}{2}\zeta'_{\Delta}(0)$, reproduces the Klebanov--Pufu--Safdi
(2011) continuum value bit-exactly for both the conformally coupled
scalar ($F_{s} = \tfrac{\log 2}{8} - \tfrac{3\zeta(3)}{16\pi^{2}}$) and
the free Camporesi--Higuchi Dirac ($F_{D} = \tfrac{\log 2}{4} +
\tfrac{3\zeta(3)}{8\pi^{2}}$); \emph{(B)} the wedge KMS entropy
$S(\rhoW) \sim 2\log\nmax$ extracts the wedge boundary dimension
$\dim(\partial W) = 2$ as a degeneracy log on the lowest $\Kalpha$
shell; \emph{(C)} the framework's modular Hamiltonian $\Kalpha =
J_{\mathrm{polar}}$ identifies, at finite cutoff, with the continuum
Casini--Huerta--Myers hemisphere modular Hamiltonian, with
GH-convergence rate inherited from~\cite{paper38, paper45}. The
scalar/Dirac partition functions decompose orthogonally under the
master Mellin engine of~\cite{paper18}: $F_{D} + 2 F_{s} = \log(2)/2$
isolates the M2 (Seeley--DeWitt) sector, $F_{D} - 2 F_{s} =
3\zeta(3)/(4\pi^{2})$ isolates the M3 (half-integer Hurwitz / odd-zeta)
sector. This is the first verification of the master Mellin engine on
non-GeoVac-internal physics observables. The bulk-side identifications
(Ryu--Takayanagi minimum surfaces, JLMS bulk-modular reconstruction,
AdS$_{4}$ / H$^{4}$ dual geometry) are blocked at the framework's
existing infrastructure: zero AdS / H$^{4}$ machinery, and the
Wick-rotation $\sthree \to H^{3}$ supplies no natural discrete
subgroup~\cite{paper29}.
\end{abstract}

\section{Introduction}
\label{sec:intro}

The GeoVac spectral triple~\cite{paper7, paper32} is a finite-cutoff
operator-system truncation of the Camporesi--Higuchi spectral
triple~\cite{camporesi_higuchi1996} on the round $\sthree$. Its
modular structure on the hemispheric wedge has been developed across
Papers~42--49~\cite{paper42, paper43, paper44, paper45, paper46,
paper47, paper48, paper49}: the four-witness Wick-rotation theorem
closes at the operator-system level (Riemannian, finite cutoff) via
the geometric BW-$\alpha$ generator $\Kalpha = J_{\mathrm{polar}}$
with integer spectrum, and via the Tomita--Takesaki BW-$\gamma$
modular operator on the GNS Hilbert--Schmidt space; the modular flow
$\sigma_{2\pi}(\Op) = \Op$ closes bit-exactly. The K$^{+}$-weak-form
Lorentzian propinquity convergence~\cite{paper45} extends the
construction to signature $(3,1)$, with the strong-form
extension~\cite{paper46} and the norm-resolvent
non-compact-carrier closure~\cite{paper47} completing the Lorentzian
arc. The Krein--Mondino--S\"amann bridge~\cite{paper48} and the
OSLPLS strong-form bridge~\cite{paper49} provide the synthetic-Lorentzian
geometric interpretation.

The present paper turns from the math.OA structural content of the
modular machinery to its \emph{physics-side observables}. We ask
whether the framework's wedge KMS state $\rhoW$, evaluated on
boundary observables corresponding to standard CFT$_{3}$-on-$\sthree$
calculations, reproduces the known continuum values, and whether the
resulting structural pattern exhibits any genuinely new content
beyond a numerical match.

\paragraph{Headline results.}
We establish three complementary boundary-side observables and one
structural decomposition theorem:

\begin{enumerate}
\item \textbf{Partition function match} (Sec.~\ref{sec:partition}):
the universal F-theorem coefficient extracted from the framework's
spectral-zeta machinery agrees \emph{bit-exactly} with the
Klebanov--Pufu--Safdi (2011)~\cite{klebanov_pufu_safdi2011} continuum
values for both the conformally coupled scalar and the free
Camporesi--Higuchi Dirac on the round $\sthree$.

\item \textbf{Master Mellin engine decomposition}
(Theorem~\ref{thm:mellin_decomp}): the scalar and Dirac partition
functions decompose into linear combinations of $\log 2$ and
$\zeta(3)/\pi^{2}$ that project orthogonally onto the M2 and M3
sectors of the master Mellin engine~\cite{paper18}. This is the
first verification of the master Mellin engine on non-GeoVac-internal
physics observables.

\item \textbf{Wedge KMS entropy structural classification}
(Sec.~\ref{sec:wedge_kms}): the von Neumann entropy of $\rhoW$ scales
as $S(\rhoW) \sim 2 \log \nmax$, with the slope $2 = \dim(\partial
\Wedge)$ identifying the wedge boundary dimension as a degeneracy
log on the lowest $\Kalpha$ shell. The F-coefficient does \emph{not}
appear in $S(\rhoW)$, demonstrating that the framework's spectral-zeta
side and wedge KMS state side encode \emph{different} continuum
observables on the same modular state.

\item \textbf{Casini--Huerta--Myers identification}
(Sec.~\ref{sec:chm}): the framework's modular Hamiltonian $\Kalpha$
identifies, at finite cutoff, with the continuum CHM hemisphere
modular Hamiltonian $K_{\mathrm{cont}} = 2\pi \int_{H_{+}} \xi \cdot
T^{00}\, d\Omega$~\cite{casini_huerta_myers2011}; convergence
$\Kalpha \to K_{\mathrm{cont}}$ as $\nmax \to \infty$ is bounded by
the Paper~38 GH-convergence rate $\Lambda \le C_{3} \cdot
\gamma_{\nmax}$ with $\gamma_{\nmax} \sim (4/\pi)\log(\nmax)/\nmax$.
\end{enumerate}

\paragraph{Bulk-side scope.}
Bulk-side identifications --- Ryu--Takayanagi minimum-surface
entanglement entropy~\cite{ryu_takayanagi2006}, the
JLMS~\cite{jafferis_lewkowycz_maldacena_suh2016} bulk-modular
reconstruction, and AdS$_{4}$ / H$^{4}$ dual geometry --- are
blocked at the framework's existing infrastructure: zero AdS$_{4}$ /
H$^{4}$ machinery exists, and a previous attempt at the natural Wick
continuation $\sthree \to H^{3}$~\cite{paper29} found no discrete
subgroup $\Gamma \subset SO(3,1)$ selectable from framework
invariants. We document this scope statement honestly in
Sec.~\ref{sec:bulk_blocked}, with no claim that the bulk side is
reachable from the framework's current axiom set.

\section{Setup}
\label{sec:setup}

\subsection{The GeoVac spectral triple on $\sthree$}

The Fock projection~\cite{fock1935, paper7} maps the
single-particle hydrogen Hamiltonian conformally to the
Laplace--Beltrami operator on $\sthree$ via stereographic projection
in momentum space. The framework's spectra on the unit $\sthree$ are
exact in $\Q$:

\begin{itemize}
\item \textbf{Scalar Laplacian.} Eigenvalues $\lambda_{n}^{2} =
n(n+2)$ with multiplicity $d_{n} = (n+1)^{2}$ for $n = 0, 1, 2,
\dots$. After the conformal mass shift $\xi R_{\mathrm{scalar}} =
\tfrac{1}{8} \cdot 6 = \tfrac{3}{4}$ for the conformally coupled
scalar, the spectrum becomes $(n + \tfrac{1}{2})(n + \tfrac{3}{2})$
with the same multiplicity.

\item \textbf{Camporesi--Higuchi Dirac.}
Eigenvalues $|\lambda_{n}| = n + \tfrac{3}{2}$ with multiplicity
$g_{n}^{\mathrm{Dirac}} = 2(n+1)(n+2)$ for the Weyl (single
chirality) sector, $4(n+1)(n+2)$ for the full $4$-component
Dirac. Both sectors carry the same Hurwitz-zeta structure under
analytic continuation~\cite{paper22, paper23}.
\end{itemize}

The truncation at level $\nmax$ retains all $n \le \nmax$ shells and
defines the operator system
$\Op_{\nmax} = P_{\nmax} \Acal P_{\nmax}$ in the
Connes--van~Suijlekom framework~\cite{connes_vs2021}. Convergence
$\Op_{\nmax} \to \Op_{\sthree}$ in the Latr\'emoli\`ere propinquity
holds at rate $\gamma_{\nmax} \sim (4/\pi) \log(\nmax) /
\nmax$~\cite{paper38}; the rate constant $4/\pi$ is the M1
Hopf-base-measure signature of the master Mellin
engine~\cite{paper18}.

\subsection{Wedge KMS state}

The hemispheric wedge $W$ corresponds, in the truncated full-Dirac
basis, to the half-space with $m_{j} > 0$ along the
wedge-fixing axis. Per~\cite{paper42} Def~5.1, the BW-$\alpha$
modular Hamiltonian on the wedge is
\[
\Kalpha = J_{\mathrm{polar}}
\]
where $J_{\mathrm{polar}}$ is the rotation generator preserving the
wedge, with integer eigenvalues $2m_{j} \in \{1, 3, 5, \dots,
2\nmax - 1\}$ and multiplicity $g(2k + 1) = (\nmax - k)(\nmax - k +
1)$ for $k = 0, \dots, \nmax - 1$. The wedge KMS state at canonical
BW normalization is
\[
\rhoW = \frac{e^{-\Kalpha}}{Z}, \qquad Z = \Tr_{\Hilb_{W}}\bigl(e^{-\Kalpha}\bigr).
\]
This is the substrate for all three boundary-side observables that
follow.

\section{Partition function: bit-exact match to KPS}
\label{sec:partition}

\subsection{Continuum reference}

For the round unit $\sthree$, the free energies of free CFTs computed
via Hurwitz-zeta regularization
are~\cite{klebanov_pufu_safdi2011}:
\begin{align}
F_{s}^{\KPS} &= \frac{\log 2}{8} - \frac{3\zeta(3)}{16\pi^{2}}
\approx 0.063807, \label{eq:Fs_KPS}\\
F_{D}^{\KPS} &= \frac{\log 2}{4} + \frac{3\zeta(3)}{8\pi^{2}}
\approx 0.218960, \label{eq:FD_KPS}
\end{align}
for the conformally coupled scalar and the free massless Weyl
Dirac, respectively. Both values are expressed in terms of the
functional determinant identities $F = -\tfrac{1}{2}\zeta'_{\Delta}(0)$
(scalar) and $F = \zeta'_{|D|}(0)$ (Dirac), with the analytic
continuation of the divergent spectral sum supplied by the standard
zeta-function regularization. The transcendental signature involves
only $\log 2$, $\zeta(3)$, and $\pi^{2}$.

\subsection{Framework's Hurwitz machinery}

The framework's spectral zeta on the truncated $\sthree$ admits a
closed-form Hurwitz expansion. For the Dirac case, the Dirichlet
series can be written exactly via the half-integer Hurwitz identity
$\zeta_{H}(s, 3/2) = (2^{s} - 1)\zeta_{R}(s) - 2^{s}$:
\begin{equation}
D_{\mathrm{Dirac}}(s) = \sum_{n = 0}^{\infty} \frac{g_{n}^{\mathrm{Dirac}}}{|\lambda_{n}|^{s}}
= 2\bigl(2^{s - 2} - 1\bigr) \zeta_{R}(s - 2) - \tfrac{1}{2}\bigl(2^{s} - 1\bigr) \zeta_{R}(s).
\label{eq:DDirac_closed_form}
\end{equation}
This identity is encoded in the production module
\texttt{geovac/qed\_two\_loop.py} and is verified symbolically across
integer $s$.

For the conformally coupled scalar, the spectral zeta admits the
Hurwitz-asymptotic expansion
\begin{equation}
\zeta_{\Delta}(s) = \sum_{n = 0}^{\infty} \frac{(n + 1)^{2}}{[(n + 1/2)(n + 3/2)]^{s}}
= \sum_{k \ge 0} g_{k}(s)\, \zeta_{R}(2s + 2k - 2),
\label{eq:zeta_scalar_expansion}
\end{equation}
with $g_{k}(s) = \tbinom{s + k - 1}{k} / 4^{k}$ the
binomial-like coefficient.

\subsection{Bit-exact match: scalar}
\label{subsec:scalar_match}

Differentiating~(\ref{eq:zeta_scalar_expansion}) at $s = 0$ and using
$g_{k}(0) = \delta_{k,0}$, $g_{k}'(0) = 1/(k \cdot 4^{k})$ for $k
\ge 1$, plus the standard identities $\zeta_{R}'(-2) = -\zeta(3)/(4
\pi^{2})$, $\zeta_{R}(0) = -\tfrac{1}{2}$, $\zeta_{R}(2k - 2)$
expressible in $\pi^{2k - 2} \cdot \Q$ for $k \ge 2$:
\begin{equation}
\zeta'_{\Delta}(0) = 2\zeta'_{R}(-2) + \sum_{k \ge 1} \frac{\zeta_{R}(2k - 2)}{k \cdot 4^{k}}.
\label{eq:zeta_prime_scalar}
\end{equation}

\begin{theorem}[Bit-exact scalar partition function]
\label{thm:scalar_bitexact}
The framework's scalar zeta-derivative satisfies
\begin{equation}
F_{s} := -\tfrac{1}{2}\,\zeta'_{\Delta}(0) = \frac{\log 2}{8} - \frac{3\zeta(3)}{16\pi^{2}} = F_{s}^{\KPS},
\end{equation}
in bit-exact agreement with~(\ref{eq:Fs_KPS}).
\end{theorem}

\begin{proof}[Proof sketch]
The asymptotic Hurwitz expansion~(\ref{eq:zeta_prime_scalar})
converges geometrically in $k$ at rate $1/4$ per term;
truncation at $k_{\max} = 100$ gives 61+ matching digits at
$\mathrm{dps} = 200$ when compared to $-2 F_{s}^{\KPS}$ via
\texttt{mpmath}. PSLQ at coefficient bound $10^{6}$ and tolerance
$10^{-40}$ finds the integer relation $[8, 2, -3]$:
\[
8 \cdot \zeta'_{\Delta}(0) + 2 \log 2 - 3 \cdot \frac{\zeta(3)}{\pi^{2}} = 0,
\]
which is exactly the claimed identity. See driver
\texttt{debug/ads\_track\_a\_scalar\_partition\_function.py} and JSON
\texttt{debug/data/ads\_track\_a\_scalar\_partition\_function.json}
for full computation.
\end{proof}

\subsection{Bit-exact match: Dirac}
\label{subsec:dirac_match}

Differentiating~(\ref{eq:DDirac_closed_form}) symbolically at $s = 0$,
using the same Riemann-zeta identities at $s = -2, 0$:
\begin{align}
D'_{\mathrm{Dirac}}(0)
&= \tfrac{1}{2}\log 2 \cdot \zeta_{R}(-2) + \bigl(-\tfrac{3}{2}\bigr) \zeta_{R}'(-2)
- \tfrac{1}{2}\bigl[\log 2 \cdot \zeta_{R}(0) + 0\bigr] \\
&= 0 - \tfrac{3}{2} \cdot \bigl(-\tfrac{\zeta(3)}{4\pi^{2}}\bigr)
- \tfrac{1}{2} \cdot \bigl(-\tfrac{\log 2}{2}\bigr)
= \frac{3\zeta(3)}{8\pi^{2}} + \frac{\log 2}{4}.
\end{align}

\begin{theorem}[Bit-exact Dirac partition function]
\label{thm:dirac_bitexact}
The framework's Dirac zeta-derivative satisfies, as an exact
symbolic identity,
\begin{equation}
F_{D} := D'_{\mathrm{Dirac}}(0) = \frac{\log 2}{4} + \frac{3\zeta(3)}{8\pi^{2}} = F_{D}^{\KPS},
\end{equation}
in bit-exact agreement with~(\ref{eq:FD_KPS}).
\end{theorem}

\begin{proof}
Direct symbolic computation; \texttt{sympy.simplify(D'(0) - F\_D)}
returns $0$. See driver
\texttt{debug/ads\_track\_a\_dirac\_partition\_function.py}.
\end{proof}

\subsection{Master Mellin engine decomposition}

The master Mellin engine of~\cite{paper18}~\S III.7 classifies the
transcendental signatures of GeoVac observables via the three
sub-mechanisms
\[
\mathcal{M}[\Tr(D^{k} \cdot e^{-t D^{2}})], \qquad k \in \{0, 1, 2\},
\]
labeled M1 (Hopf-base measure, $\sqrt{\pi} \cdot \Q$ ring at $k = 0$),
M2 (Seeley--DeWitt heat-kernel, $\sqrt{\pi} \cdot \Q \oplus \pi^{2}
\cdot \Q$ ring at $k = 2$), and M3 (half-integer Hurwitz /
vertex-parity, $\zeta(3)/\pi^{2}$ and similar odd-zeta-over-pi
ratios at $k = 1$).

Theorems~\ref{thm:scalar_bitexact} and~\ref{thm:dirac_bitexact}
exhibit $F_{s}$ and $F_{D}$ as linear combinations of $\log 2$ (an
M2 transcendental, present via $\zeta_{R}'(0) = -\tfrac{1}{2}
\log(2\pi)$ in the underlying heat-kernel asymptotic) and
$\zeta(3)/\pi^{2}$ (an M3 transcendental, present via
$\zeta_{R}'(-2) = -\zeta(3)/(4\pi^{2})$ in the half-integer
Hurwitz-zeta derivative).

\begin{theorem}[Scalar/Dirac M2--M3 orthogonal decomposition]
\label{thm:mellin_decomp}
The free energies $F_{s}$ and $F_{D}$ project orthogonally onto the
M2 and M3 axes of the master Mellin engine via the dual-basis
combinations
\begin{align}
F_{D} + 2\,F_{s} &= \frac{\log 2}{2}
\qquad\text{(isolates M2; all $\zeta(3)$ cancels),}
\label{eq:M2_isolation}\\
F_{D} - 2\,F_{s} &= \frac{3\zeta(3)}{4\pi^{2}}
\qquad\text{(isolates M3; all $\log 2$ cancels).}
\label{eq:M3_isolation}
\end{align}
\end{theorem}

\begin{proof}
Substituting~(\ref{eq:Fs_KPS}) and~(\ref{eq:FD_KPS}):
\[
F_{D} + 2 F_{s} = \bigl(\tfrac{\log 2}{4} + \tfrac{3\zeta(3)}{8\pi^{2}}\bigr)
+ 2 \bigl(\tfrac{\log 2}{8} - \tfrac{3\zeta(3)}{16\pi^{2}}\bigr)
= \tfrac{\log 2}{4} + \tfrac{\log 2}{4} = \tfrac{\log 2}{2},
\]
and similarly $F_{D} - 2 F_{s} = 3 \zeta(3)/(4\pi^{2})$. The two
combinations span a $2$-dimensional space; the projection
coefficients $(1, +2)$ and $(1, -2)$ act as a dual basis for the
M2/M3 decomposition on the (scalar, Dirac) plane.
\end{proof}

\begin{remark}[First verification on non-GeoVac-internal data]
\label{rem:first_external_verification}
Theorem~\ref{thm:mellin_decomp} is the first verification of the
master Mellin engine on observables that are \emph{not}
GeoVac-internal. Previous verifications --- the case-exhaustion
theorem of~\cite{paper32}~\S VIII, the 208/208 Prediction~1 panel
of~\cite{paper35}, the modular-residual closed form of Sprint MR-B
--- all operated on framework-internal quantities. The KPS
partition functions~(\ref{eq:Fs_KPS}),~(\ref{eq:FD_KPS}) are
independently-published continuum CFT$_{3}$ data, and the framework's
spectral-zeta machinery reproduces them with the exact M2/M3
transcendental signature the master Mellin engine predicts. This
elevates the engine from a structural classification of
GeoVac-internal observables to a predictive classification of CFT
physics on the round $\sthree$ at finite cutoff.
\end{remark}

\section{Wedge KMS entropy: state-side observable}
\label{sec:wedge_kms}

\subsection{The framework's wedge KMS entropy}

The von Neumann entropy of the wedge KMS state $\rhoW$ at canonical
BW normalization $\beta = 2\pi$ is computed directly by
diagonalizing $\Kalpha$ and applying $S(\rhoW) = -\Tr[\rhoW \log
\rhoW]$. Sprint BH-Phase0
(2026-05-22)~\cite{debug:bh_phase0} first reported a small-cutoff
log-linear fit $S(\rhoW) \approx 1.94 \log\nmax + 0.56$ over $\nmax
\in \{2, \dots, 7\}$. A subsequent re-analysis (forcing-catalogue
Door~2, 2026-06-01, \texttt{debug/door2\_bw\_entropy\_reread\_memo.md})
replaced this fit by the exact $\nmax \to \infty$ resummation of the
closed-form entropy $S = \log Z + \langle \Kalpha \rangle$ (which
reproduces the direct diagonalization to $\le 4 \times 10^{-16}$).
The leading $\nmax^{2}$ shell degeneracy $g(2k+1) = (\nmax - k)(\nmax
- k + 1)$ forces the slope to be \emph{exactly} $2$, and the additive
constant has the exact closed form
\begin{equation}
S(\rhoW) = 2 \log \nmax + \big(\coth 1 - \log(2 \sinh 1)\big)
+ O(1/\nmax), \qquad \coth 1 - \log(2 \sinh 1) = 0.458448\ldots,
\label{eq:S_W_fit}
\end{equation}
the constant verified by PSLQ to residual $< 10^{-80}$ (integer
relation $[-1, 1, -1, -1]$ against $\{C_\infty, \coth 1, \log 2, \log
\sinh 1\}$). The windowed least-squares slope climbs monotonically
$1.94 \to 2.00005$ as the fit window moves to $\nmax \in [350,
2000]$; the original $1.94$ is a small-$\nmax$ artifact of a sequence
converging to exactly $2$. The slope $2 = \dim(\partial W) =
\dim(S^{2}_{\mathrm{equator}})$ identifies the wedge boundary
dimension as a degeneracy log on the lowest $\Kalpha$ shell. The
additive constant lives in the Boltzmann-base ring (functions of $e$
via $\coth 1$, $\sinh 1$, tied to the integer spacing
$\Delta\Kalpha = 2$ at $\beta = 2\pi$); it carries no master Mellin
engine transcendental, consistent with
Prop.~\ref{prop:F_vs_S_decomposition}. A PSLQ search against the
F-theorem ring $\{1, \log 2, \zeta(3)/\pi^{2}\}$ at coefficient
ceiling $10^{8}$ returns null: the wedge-entropy constant and the
F-coefficient lie in disjoint rings, sharpening
Prop.~\ref{prop:F_vs_S_decomposition} from ``F absent from $S(\rhoW)$''
to ``the state-side constant is provably outside the spectral-zeta
ring.''

\subsection{Mechanism: degeneracy log on the lowest $\Kalpha$ shell}

The $\Kalpha$ spectrum is bounded above by $2\nmax - 1$ at finite
cutoff. At BW canonical normalization, the per-state weight at
$\Kalpha = 1$ (the equator shell) is $e^{-1}/Z$, while at $\Kalpha = 3$
it drops to $e^{-3}/Z$, a factor of $e^{-2} \approx 0.135$ suppression
per shell. The partition function is therefore dominated by the
equator shell of multiplicity $n_{\mathrm{eq}} = \nmax(\nmax + 1)
\sim \nmax^{2}$:
\[
Z \approx n_{\mathrm{eq}} \cdot e^{-1} \approx 0.368\, \nmax^{2},
\]
giving effective probability $p_{\mathrm{eq}} \approx 1/n_{\mathrm{eq}}$
per equator state, and hence
\[
S(\rhoW) \approx -n_{\mathrm{eq}} \cdot p_{\mathrm{eq}} \log p_{\mathrm{eq}}
= \log n_{\mathrm{eq}} \approx 2 \log \nmax.
\]
The two asymptotic limits in $s$-deformation $\rho(s) = e^{-s
\Kalpha}/Z$, $s \to 0^{+}$ (maximally mixed on $\Hilb_{W}$, giving
$S \to \log \dim \Hilb_{W}$) and $s \to \infty$ (maximally mixed on
the equator shell, giving $S \to \log n_{\mathrm{eq}}$), are
saturated bit-exactly~\cite{debug:bh_phase0}.

\subsection{Structural decomposition: F is in the spectral-zeta side, not the state side}

\begin{proposition}[F vs $S(\rhoW)$ live in different sectors]
\label{prop:F_vs_S_decomposition}
The framework's wedge KMS state $\rhoW$ admits two structurally
distinct calculational operations that extract different continuum
observables:

\begin{itemize}
\item \textbf{Spectral-zeta side:} the analytic continuation
$-\tfrac{1}{2} \zeta'_{\Delta}(0)$ extracts the universal F-theorem
coefficient $F_{s}^{\KPS}$ via Theorem~\ref{thm:scalar_bitexact};
similarly for $F_{D}^{\KPS}$ via Theorem~\ref{thm:dirac_bitexact}.
The transcendental signature is in the master Mellin engine M2/M3
ring $\Q \cdot \log 2 \oplus \Q \cdot \zeta(3)/\pi^{2}$.

\item \textbf{Wedge KMS state side:} the von Neumann entropy
$S(\rhoW) = -\Tr[\rhoW \log \rhoW]$ extracts the wedge boundary
dimension $\dim(\partial W)$ via degeneracy log on the lowest
$\Kalpha$ shell, with leading scaling $\dim(\partial W) \cdot \log
\nmax$. The result is ring-free (a logarithm of an integer).
\end{itemize}

The two operations act on the same modular state $\rhoW$ but extract
\emph{categorically different} continuum observables: the universal
F-coefficient versus the wedge boundary dimension. The F-coefficient
does \emph{not} appear in $S(\rhoW)$.
\end{proposition}

This decomposition is consistent with the PSLQ-disjointness finding
of Sprint TD Track~5 (2026-05-09)~\cite{debug:sprint_td_track5} that
GeoVac entropies generally live outside the master Mellin engine
ring, while spectral-side observables generally live inside it.

\subsection{Three blockers for literal Ryu--Takayanagi on the framework's graph}

The continuum Ryu--Takayanagi formula~\cite{ryu_takayanagi2006}
identifies the boundary entanglement entropy of a region $A$ with
the area of the bulk minimum surface anchored at $\partial A$,
weighted by $1/(4 G_{N})$. A literal discrete analog on the
framework's graph would require: (i) a well-defined notion of
graph minimum-cut between spatial sub-regions; (ii) a perfect-tensor
structure at graph vertices ensuring isometric bulk-boundary
encoding; (iii) a bulk dual geometry.

\begin{proposition}[Three blockers for literal RT on the framework]
\label{prop:RT_blocked}
The framework's discrete $\sthree$ graph at finite $\nmax$ exhibits
three independent structural obstructions to a literal RT formula:
\begin{enumerate}
\item The Fock graph at level $\nmax$ has $\beta_{0} = \nmax$
connected components (one per $\ell$-sector), so the naive
graph-cut between spatial sub-regions is identically zero or
trivially confined to a single $\ell$-sector.

\item The Wigner 3j symbols at graph vertices are Clebsch--Gordan
projections rather than isometries; the HaPPY tensor-network
RT proof~\cite{pastawski_yoshida_harlow_preskill2015} requires
perfect tensors at network nodes.

\item No bulk-dual construction exists in the framework's
infrastructure. The natural candidate --- Wick continuation
$\sthree \to H^{3}$ --- was closed by Sprint RH-B
(2026-04-17)~\cite{paper29} as a clean structural dead end: no
discrete subgroup $\Gamma \subset SO(3,1)$ is selectable from
framework invariants.
\end{enumerate}
\end{proposition}

These three blockers are not in tension with the boundary-side
results of Sec.~\ref{sec:partition} or Sec.~\ref{sec:chm}: the
framework reproduces boundary observables on the same wedge KMS
state, but cannot independently construct a bulk dual against which
to compare via RT.

\section{Casini--Huerta--Myers identification}
\label{sec:chm}

\subsection{Continuum CHM hemisphere modular Hamiltonian}

The Casini--Huerta--Myers
formula~\cite{casini_huerta_myers2011} for the continuum modular
Hamiltonian of a hemisphere of $\sthree$ in free CFT$_{3}$ is
\begin{equation}
K_{\mathrm{cont}} = 2\pi \int_{H_{+}} \xi(x) \cdot T^{00}(x)\, d\Omega(x),
\label{eq:K_chm}
\end{equation}
where $\xi$ is the wedge-preserving Killing vector field of the
rotation that fixes the equator $S^{2}_{\mathrm{equator}}$, and
$T^{00}$ is the local energy density. The flow generated by
$K_{\mathrm{cont}}/(2\pi)$ is $2\pi$-periodic, reflecting the
canonical Bisognano--Wichmann period of the modular
flow~\cite{bisognano_wichmann1976}; the eigenvalues of
$K_{\mathrm{cont}}$ in the single-particle basis are $2\pi J$ for
integer $J = 1, 2, 3, \dots$.

\subsection{Framework identification}

The framework's wedge modular Hamiltonian (Paper~42 Def~5.1) is
\[
\Kalpha = J_{\mathrm{polar}},
\]
the rotation generator of the wedge-preserving $SO(2)$ on the
truncated Camporesi--Higuchi spectral triple. Its spectrum is
\[
\mathrm{spec}(\Kalpha) = \{2 m_{j} : m_{j} = 1/2, 3/2, \dots, (2\nmax - 1)/2\}
= \{1, 3, 5, \dots, 2\nmax - 1\},
\]
integer-valued in the canonical $(2\pi)^{-1}$ rescaling, with
multiplicity $g(2k + 1) = (\nmax - k)(\nmax - k + 1)$. The
modular flow $\sigma_{2\pi}(\Op) = \Op$ closes
bit-exactly at every tested $\nmax$~\cite{paper42}.

\begin{theorem}[CHM identification at finite cutoff]
\label{thm:chm_identification}
The framework's modular Hamiltonian $\Kalpha$ at finite $\nmax$ is
the operator-system-level analog of the continuum CHM modular
Hamiltonian $K_{\mathrm{cont}}$ of~(\ref{eq:K_chm}), with:

\begin{itemize}
\item Both spectra integer-valued in canonical $(2\pi)^{-1}$ units.
\item Both generate $2\pi$-periodic modular flow.
\item The framework's odd-integer spectrum $\{1, 3, \dots, 2\nmax - 1\}$
corresponds to the half-integer Casimir grading of the
wedge-preserving Killing field on the Weyl spinor sector.
\item Convergence $\Kalpha \to K_{\mathrm{cont}}$ as $\nmax \to
\infty$ holds in the Latr\'emoli\`ere propinquity at rate
\[
\Lambda(\Tcal_{\nmax}, \Tcal_{\sthree}) \le C_{3} \cdot \gamma_{\nmax}
\]
with $C_{3} \to 1$ and $\gamma_{\nmax} \sim (4/\pi) \log(\nmax) /
\nmax$ inherited from Paper~38~\cite{paper38} via Paper~45~\cite{paper45}.
\end{itemize}
\end{theorem}

\begin{proof}[Proof sketch]
The wedge-preserving Killing field of the round $\sthree$ rotates
the spinor harmonics by their $m_{j}$ eigenvalue; the canonical
$2\pi$-period of the rotation gives the integer $2 m_{j}$ spectrum
in the canonical normalization. The framework's truncation
$\Kalpha = J_{\mathrm{polar}}$ acts on the projected wedge
sub-space $P_{\Wedge} \Hilb_{\mathrm{GV}}^{\nmax}$ with exactly
the same $2 m_{j}$ eigenvalues, restricted to the multiplicities
allowed by the truncation. Convergence in the propinquity follows
from~\cite{paper38, paper45}.
\end{proof}

\section{Bulk-side scope}
\label{sec:bulk_blocked}

The three boundary-side observables (Sec.~\ref{sec:partition},
Sec.~\ref{sec:wedge_kms}, Sec.~\ref{sec:chm}) all reproduce or
identify their continuum CFT$_{3}$-on-$\sthree$ analogs at finite
cutoff. The corresponding \emph{bulk-side} identifications --- AdS$_{4}$
gravity, JLMS bulk-modular reconstruction, RT minimum surfaces,
holographic entanglement entropy --- are blocked at the framework's
existing infrastructure.

\begin{proposition}[Bulk side blocked]
\label{prop:bulk_blocked}
None of the bulk-side identifications corresponding to
Sec.~\ref{sec:partition}--\ref{sec:chm} are reachable from the
framework's current infrastructure:

\begin{enumerate}
\item No AdS$_{4}$ / H$^{4}$ machinery exists in the framework's
codebase (\texttt{geovac/}); a search for AdS, hyperbolic-4-space,
or holographic primitives returns no matches.

\item Sprint RH-B (2026-04-17)~\cite{paper29} closed the natural
candidate --- Wick continuation $\sthree \to H^{3}$ --- as a clean
structural dead end: the continuation supplies no natural discrete
subgroup $\Gamma \subset SO(3,1)$ selectable from framework
invariants.

\item Sprint L3e-P3 (2026-05-23) found that Paper~38's $4/\pi$ rate
does \emph{not} transport to the non-compact Coulomb setting,
suggesting that any discrete-to-bulk convergence statement would
require a non-compact rate analog that is itself an open NCG-math
problem.
\end{enumerate}
\end{proposition}

The structural reading of Prop.~\ref{prop:bulk_blocked} is that the
bulk-side identification falls into Class~1 (calibration-external)
of the framework's external-input partition~\cite{debug:external_input_three_class_partition}:
the framework determines the boundary-side observables structurally
but does not autonomously generate the bulk-dual geometry. Importing
AdS$_{4}$ infrastructure would be a multi-month engineering
commitment outside the present paper's scope.

\section{Extension to $S^{5}$: structural refinement of the dual-basis decomposition}
\label{sec:s5_extension}

The bit-exact match of Sec.~\ref{sec:partition} extends to round $S^{5}$
via the same Hurwitz-zeta machinery. The Camporesi--Higuchi spectra on
$S^{5}$ are spectrally analogous to $S^{3}$: the conformally coupled
scalar has eigenvalues $(n + 3/2)(n + 5/2)$ with multiplicity
$(2n + 4)(n + 1)(n + 2)(n + 3)/24$ for $n = 0, 1, 2, \dots$, and the
Weyl Camporesi--Higuchi Dirac has $|\lambda_{n}| = n + 5/2$ with
multiplicity $(n + 1)(n + 2)(n + 3)(n + 4)/12$. The framework's
spectral-zeta machinery on these spectra produces partition function
values in a closed-form ring extending~\cite{klebanov_pufu_safdi2011}'s
S$^{3}$ ring by one $\zeta(5)/\pi^{4}$ generator.

\subsection{Bit-exact closed forms on $S^{5}$}

\begin{theorem}[Conformally coupled scalar on $S^{5}$]
\label{thm:scalar_S5}
The framework's scalar zeta-derivative on round unit $S^{5}$ satisfies
\begin{equation}
F_{s}^{S^{5}} := -\tfrac{1}{2}\,\zeta'_{\Delta_\text{conf}}(0)^{S^{5}}
= -\frac{\log 2}{128} - \frac{\zeta(3)}{128\pi^{2}} + \frac{15\,\zeta(5)}{256\pi^{4}}
\approx -5.74 \times 10^{-3}.
\label{eq:Fs_S5}
\end{equation}
\end{theorem}

\begin{proof}[Proof sketch]
The Hurwitz expansion
\begin{equation}
\zeta_{S^{5}}^\text{conf}(s) = \frac{1}{12}\sum_{k \ge 0} g_{k}(s) [\zeta_{R}(2s + 2k - 4) - \zeta_{R}(2s + 2k - 2)]
\end{equation}
with $g_{k}(s) = \binom{s + k - 1}{k}/4^{k}$ follows from the
re-indexed multiplicity structure $v^{2}(v^{2} - 1)/12$ at $v = n + 2$.
At $s = 0$, only $g_{0}(0) = 1$ survives, giving
$\zeta_{S^{5}}^\text{conf}(0) = 0$; the derivative
\begin{equation}
\zeta'_{S^{5}}(0) = \frac{1}{6}[\zeta_{R}'(-4) - \zeta_{R}'(-2)]
+ \frac{1}{12}\sum_{k \ge 1} \frac{\zeta_{R}(2k - 4) - \zeta_{R}(2k - 2)}{k \cdot 4^{k}}
\end{equation}
converges geometrically in $k$. PSLQ at 200~dps and $k_{\max} = 100$
finds the integer relation $[32, -2, -2, 15]$ (applied to
$4\,\zeta'_{S^{5}}(0)$) bit-exactly identifying
\begin{equation}
\zeta'_{S^{5}}(0) = \frac{\log 2}{64} + \frac{\zeta(3)}{64\pi^{2}}
- \frac{15\,\zeta(5)}{128\pi^{4}}.
\end{equation}
The standard identities $\zeta_{R}'(-2) = -\zeta(3)/(4\pi^{2})$ and
$\zeta_{R}'(-4) = +3\zeta(5)/(4\pi^{4})$ deliver the $\zeta(5)/\pi^{4}$
contribution; the $\log 2$ piece arises from
$\zeta_{R}'(0) = -\tfrac{1}{2}\log(2\pi)$ via the Hurwitz
analytic-continuation tail.
\end{proof}

\begin{theorem}[Weyl Camporesi--Higuchi Dirac on $S^{5}$]
\label{thm:dirac_S5}
The framework's Weyl Camporesi--Higuchi Dirac zeta-derivative on round
unit $S^{5}$ satisfies, as an exact symbolic identity,
\begin{equation}
D'_\text{Dirac}(0)^{S^{5}}
= -\frac{3\log 2}{128} - \frac{5\,\zeta(3)}{128\pi^{2}} - \frac{15\,\zeta(5)}{256\pi^{4}}
\approx -0.02163.
\label{eq:DD_S5}
\end{equation}
\end{theorem}

\begin{proof}[Proof sketch]
The Dirichlet series
\begin{equation}
D_\text{Dirac}(s)^{S^{5}} = \frac{1}{12} \left[\zeta_{H}(s - 4, 5/2) - \tfrac{5}{2}\zeta_{H}(s - 2, 5/2) + \tfrac{9}{16}\zeta_{H}(s, 5/2)\right]
\label{eq:DD_S5_closed}
\end{equation}
follows from re-indexing $u = n + 5/2$ in the multiplicity
$(n + 1)(n + 2)(n + 3)(n + 4)/12 = (u^{2} - 9/4)(u^{2} - 1/4)/12$. The
half-integer Hurwitz function $\zeta_{H}(s, 5/2)$ has closed form
$(2^{s} - 1)\zeta_{R}(s) - 2^{s} - (2/3)^{s}$. Differentiating each
shift symbolically at $s = 0$ and combining via the multiplicity
coefficients $(1, -5/2, 9/16)$ delivers the claimed result.
\end{proof}

\subsection{Log 3 cancellation theorem}

\begin{theorem}[$\log 3$ cancellation on $S^{5}$]
\label{thm:log3_cancellation}
The $(2/3)^{s}$ terms in $\zeta_{H}(s - a, 5/2)$ for
$a \in \{0, 2, 4\}$ contribute, at $s = 0$, factors of
$\log(2/3) = \log 2 - \log 3$ scaled by $(2/3)^{-a} = (3/2)^{a}$. The
coefficient of $\log 3$ in $D'_\text{Dirac}(0)^{S^{5}}$ evaluates to
\begin{equation}
\frac{1}{12}\left[\frac{81}{16} - \frac{5}{2}\cdot\frac{9}{4} + \frac{9}{16}\cdot 1\right]
= \frac{1}{12}\left[\frac{81 - 90 + 9}{16}\right] = 0
\end{equation}
exactly.
\end{theorem}

\begin{remark}
The structural combination $(1, -5/2, 9/16)$ encoded in the
multiplicity polynomial $(u^{2} - 9/4)(u^{2} - 1/4) = u^{4} - (5/2) u^{2} + 9/16$
aligns the three shift operators $\{\zeta_{H}(s, 5/2),
\zeta_{H}(s-2, 5/2), \zeta_{H}(s-4, 5/2)\}$ such that all $\log 3$
contributions cancel exactly. The framework's spectral content on
$S^{5}$ stays in the M-engine ring $\{\log 2, \zeta(3)/\pi^{2},
\zeta(5)/\pi^{4}\}$, even though the naive Hurwitz expansion suggests
a contribution from $\log 3$. The same mechanism extends the M2/M3
ring statement from $S^{3}$ to $S^{5}$.
\end{remark}

\subsection{Dual-basis projection: structural-specificity to $S^{3}$}

Theorem~\ref{thm:mellin_decomp} of Sec.~\ref{sec:partition} establishes
that on $S^{3}$, the (scalar, Dirac) pair admits a dual-basis
projection onto the M2/M3 axes: $F_{D} + 2 F_{s} = \log(2)/2$ (M2
only), $F_{D} - 2 F_{s} = 3\zeta(3)/(4\pi^{2})$ (M3 only). On $S^{5}$,
the analogous statement \emph{fails}.

\begin{proposition}[Dual-basis non-extension to $S^{5}$]
\label{prop:dual_basis_S5}
The Paper~50 Theorem~\ref{thm:mellin_decomp} dual-basis projection
does not extend to $S^{5}$. The M-engine ring on $S^{5}$ is
three-dimensional ($\log 2$, $\zeta(3)/\pi^{2}$, $\zeta(5)/\pi^{4}$),
but the (scalar, Dirac) plane is two-dimensional. The system
$a F_{s}^{S^{5}} + b F_{D}^{S^{5}} = c \cdot (\text{single M-axis})$
is over-determined: imposing isolation of one M-axis requires three
linear conditions on two parameters $(a, b)$, which has no nontrivial
solution.

Explicitly, attempting to isolate $\log 2$:
\begin{align}
\zeta(3)/\pi^{2}\text{ coefficient} = 0 \;&\Longrightarrow\; a + 5b = 0 \;\Longrightarrow\; a = -5b,\\
\zeta(5)/\pi^{4}\text{ coefficient} = 0 \;&\Longrightarrow\; 15a - 15b = 0 \;\Longrightarrow\; a = b,
\end{align}
which are inconsistent unless $b = 0$.
\end{proposition}

\begin{remark}[Continuation of the $F$-coefficient ladder through $S^{11}$; dual-basis non-extension]
\label{rem:general_d_conjecture}
Two statements must be distinguished here: the \emph{single-species}
in-ring closure of $F$ (which continues to all odd $d$ tested), and the
\emph{dual-basis projection} of Theorem~\ref{thm:mellin_decomp} (which is
special to $S^{3}$). The structural-specificity of
Theorem~\ref{thm:mellin_decomp} concerns only the second.

\textbf{Single-species in-ring closure (continues through $S^{11}$).}
For each free species individually, $F = -\tfrac{1}{2}\zeta'_{\Delta}(0)$
on round $S^{d}$ ($d$ odd) lies in the ring
\[
\mathcal{R}_{d} = \{\log 2\} \cup
\{\zeta(2k+1)/\pi^{2k} : k = 1, \dots, (d-1)/2\}
\]
of dimension $(d-1)/2 + 1$ (top atom $\zeta(d)/\pi^{d-1}$). This is
verified bit-exactly at $S^{3}$, $S^{5}$
(Theorems~\ref{thm:scalar_S5},~\ref{thm:dirac_S5}), and --- by direct
high-precision computation (Door~1 graduation, 2026-06-03; $260$ dps,
frozen basis, reconstruction error $\lesssim 10^{-261}$, decoy-controlled)
--- at $S^{7}$, $S^{9}$, and $S^{11}$. The conformally coupled scalar on
$S^{11}$, for instance, is
\begin{align}
\zeta'_{\mathrm{sc}, S^{11}}(0) =\;
&-\tfrac{7}{131072}\log 2
 -\tfrac{3897}{45875200}\tfrac{\zeta(3)}{\pi^{2}}
 -\tfrac{485}{8257536}\tfrac{\zeta(5)}{\pi^{4}} \notag\\
&+\tfrac{609}{1310720}\tfrac{\zeta(7)}{\pi^{6}}
 +\tfrac{425}{262144}\tfrac{\zeta(9)}{\pi^{8}}
 +\tfrac{1023}{524288}\tfrac{\zeta(11)}{\pi^{10}},
\end{align}
and the free Weyl Camporesi--Higuchi Dirac value is symbolic-exact in the
same ring. A decoy test assigns Catalan $G$, $\log 3$, and unscaled
$\zeta(3)$ coefficient exactly zero, confirming a clean M3 signature.

\emph{Erratum.} An earlier version of this remark reported $S^{7}$ as a
structural non-match (``no PSLQ relation in the tested bases''). That was a
\emph{false negative} from an under-resolved search. The genuine $S^{7}$
scalar relation
\[
\zeta'_{\mathrm{sc}, S^{7}}(0) =
-\tfrac{1}{512}\log 2
-\tfrac{41}{15360}\tfrac{\zeta(3)}{\pi^{2}}
+\tfrac{5}{1024}\tfrac{\zeta(5)}{\pi^{4}}
+\tfrac{63}{2048}\tfrac{\zeta(7)}{\pi^{6}}
\]
has integer coefficients of order $10^{4}$, but locking it requires the
value to $\gtrsim 200$ digits; a $30$-digit value run against
$\mathrm{maxcoeff} = 10^{16}$ sits below the resolution threshold and
returns spurious non-detection. Verified independently at $200$ dps
(reconstruction error $\approx 2.4 \times 10^{-202}$) and reproduced by the
Door~1 driver at $260$ dps. The ladder therefore \emph{generates}: every
odd rung $S^{3}, \dots, S^{11}$ has a rational-coefficient closed form in
$\mathcal{R}_{d}$, with no transcendental outside $\mathcal{R}_{d}$.

\textbf{Dual-basis projection (special to $S^{3}$, unaffected).} Isolating
a single $M$-axis using only the $(\text{scalar}, \text{Dirac})$ pair
requires $\dim \mathcal{R}_{d} = 2$, which holds at $S^{3}$ alone. For
$d \ge 5$ the projection is over-determined
(Prop.~\ref{prop:dual_basis_S5}); the related cross-species ladder
recursion
$\mathrm{scalar} - m \cdot \mathrm{Dirac} = c_{d} \cdot \mathrm{Dirac}_{S^{d-2}}$
likewise closes at $S^{5}$ only, with the top-atom ratio
$\mathrm{scalar}_{\mathrm{top}}/\mathrm{Dirac}_{\mathrm{top}} =
2^{-(d-5)/2}$ hitting the required value ($m = 1$) at $d = 5$ alone
(Door~1 graduation). Thus the $F$-coefficients are in-ring at every odd
$d$, but the two-species dual basis that \emph{spans} the ring is special
to $S^{3}$. Supplying the missing higher-spin free species ---
$(d-1)/2 - 1$ of them at $S^{d}$ --- to span $\mathcal{R}_{d}$ remains the
natural continuation, and is the actual structural-specificity content of
Theorem~\ref{thm:mellin_decomp}.
\end{remark}

\subsection{Implementation effort and continuation pattern}

The S$^{5}$ extension was computed in ~35 minutes of focused
main-session work, continuing the discrete-tractability advantage
documented across the math.OA arc. The
\texttt{debug/ads\_track\_a\_s5\_\{scalar, dirac\}\_partition\_function.py}
scripts implement the Hurwitz expansions and PSLQ verifications;
\texttt{debug/ads\_track\_a\_s5\_extension\_memo.md} documents the
findings in full.

\section{Synthesis and open questions}
\label{sec:synthesis}

\subsection{Unified picture}

The framework's wedge KMS state $\rhoW = e^{-\Kalpha}/Z$ supports
three structurally distinct calculational operations, each
extracting a different continuum CFT$_{3}$-on-$\sthree$ boundary
observable:

\begin{center}
\begin{tabular}{lll}
\toprule
Operation & Continuum analog & Ring \\
\midrule
$-\tfrac{1}{2}\zeta'_{\Delta}(0)$ (Theorem~\ref{thm:scalar_bitexact}) & F-theorem coefficient & M2 + M3 \\
$-\Tr[\rhoW \log \rhoW]$ (Prop.~\ref{prop:F_vs_S_decomposition}) & Wedge boundary dimension & ring-free \\
$\mathrm{spec}(\Kalpha)$ (Theorem~\ref{thm:chm_identification}) & CHM modular Hamiltonian & ring-free \\
\bottomrule
\end{tabular}
\end{center}

All three converge to their continuum analogs as $\nmax \to \infty$
per~\cite{paper38, paper45}.

\subsection{Open questions}

\begin{enumerate}
\item \textbf{Other free-field CFTs.} The scalar and Dirac cases
exhibit the orthogonal M2/M3 decomposition of
Theorem~\ref{thm:mellin_decomp}. Does the same projection structure
extend to higher-spin partition functions on $\sthree$, in
particular the Maxwell ($p = 1$ form) and graviton ($p = 2$) cases?
The Maxwell partition function on round $\sthree$ has known closed
forms~\cite{beccaria_tseytlin2017}; testing whether the framework's
spectral-zeta machinery reproduces it (and whether the M2/M3
decomposition extends) is a natural follow-up.

\item \textbf{Squashed-$\sthree$ generalization.} The KPS values
above are for the round (unit) $\sthree$. The squashed-$\sthree$
partition functions of supersymmetric and free
CFTs~\cite{jafferis_klebanov_pufu_safdi2011} use deformed half-integer
Hurwitz machinery that may not reduce to the framework's MR-B
Jacobi-$\vartheta_{2}$ closed form. Whether the framework can match
the squashed case is open.

\item \textbf{Propinquity-rate convergence statement.} Theorem~\ref{thm:chm_identification}
states convergence at the propinquity-bound level. A quantitative
rate for the partition-function-side convergence $|F_{\nmax} -
F^{\KPS}| \le C \cdot \gamma_{\nmax}$ with explicit $C$ requires a
Mellin-domain version of the Berezin-reconstruction framework of
Paper~38; this is open work.

\item \textbf{Bulk-side identification (Prop.~\ref{prop:bulk_blocked}).}
The JLMS bulk-modular Hamiltonian identification, the RT
entanglement-entropy formula, and the AdS$_{4}$ dual geometry
remain open. Importing AdS$_{4}$ / H$^{4}$ infrastructure is the
named multi-month engineering commitment that would reach these.
\end{enumerate}

\subsection{Catalogue and open targets}
\label{sec:catalogue}

The bit-exact matches at $S^{3}$ and $S^{5}$ exhibit a pattern: the
framework's discrete spectral-zeta machinery reproduces
\emph{independently-published} continuum CFT-on-sphere free-energy
values, and the transcendental signature decomposes via the master
Mellin engine~\cite{paper18}. This pattern admits systematic
extension. The following table catalogues the current status:

\begin{center}
\begin{tabular}{lccc}
\toprule
System & Continuum reference & Framework status & M-engine reading \\
\midrule
Free scalar (conformal) on $\sthree$ & KPS~\cite{klebanov_pufu_safdi2011} & DONE Sec.~\ref{sec:partition} & M2 + M3 \\
Free Dirac (Weyl, CH) on $\sthree$ & KPS~\cite{klebanov_pufu_safdi2011} & DONE Sec.~\ref{sec:partition} & M2 + M3 \\
Free scalar (conformal) on $S^{5}$ & Beccaria-Tseytlin~\cite{beccaria_tseytlin2017} & DONE Sec.~\ref{sec:s5_extension} & M2 + M3 extended \\
Free Dirac (Weyl, CH) on $S^{5}$ & Beccaria-Tseytlin~\cite{beccaria_tseytlin2017} & DONE Sec.~\ref{sec:s5_extension} & M2 + M3 extended \\
Free scalar (conformal) on $S^{7}$ & --- & numerical only; PSLQ fails on simple ring & UNKNOWN (Remark~\ref{rem:general_d_conjecture}) \\
\midrule
Free Maxwell on $\sthree$ & 3D scalar duality~\cite{beccaria_tseytlin2017} & open: $\sthree$ Maxwell-scalar duality gives same as scalar & expected M2 + M3 \\
Free Maxwell on $S^{5}$ & BT~\cite{beccaria_tseytlin2017} & TODO; needs transverse 1-form Hodge spectrum & open \\
Free higher-rank tensor ($p \ge 2$) on $S^{d}$ & BT~\cite{beccaria_tseytlin2017}, Anninos & TODO; needs $\mathrm{SO}(d+1)$ rep theory & open \\
Free (non-)minimal scalar (zero mode) on $S^{d}$ & --- & open: $m{=}2$ edge correction introduces $\log 3 + \log 5$ at $S^{5}$ & ring extension needed \\
Squashed $\sthree$ partition function & JKPS, Hama-Hosomichi-Lee & open: deformed half-integer Hurwitz not in MR-B Jacobi-$\vartheta_{2}$ closed form & may require new structural mechanism \\
\bottomrule
\end{tabular}
\end{center}

\paragraph{Holographic Weyl anomaly $\leftrightarrow$ Connes--Chamseddine spectral action: both live in M2.}
A structural correspondence with the AdS/CFT literature that the framework
inherits directly --- without additional computation --- is the M2-sector
identification of the holographic Weyl anomaly and the
Connes--Chamseddine spectral action.

Henningson and Skenderis~\cite{henningson_skenderis1998} derived the Weyl
anomaly of CFT$_{d}$ on a manifold $M$ via the bulk AdS$_{d+1}$
gravitational action regularized at a finite radial cutoff. For even
$d$, the result reproduces the boundary central charges
$a_{\text{CFT}}, c_{\text{CFT}}$ as specific combinations of bulk
gravitational coefficients. The bulk integral is computed via the
Riemannian curvature data of AdS, evaluated through a heat-kernel
expansion in the AdS-radius cutoff $\epsilon$.

The Connes--Chamseddine spectral action~\cite{chamseddine_connes1997}
$\mathrm{Tr}\,f(D/\Lambda)$, on any almost-commutative spectral triple,
has heat-kernel expansion coefficients given by Seeley--DeWitt
coefficients $a_{k}(D^{2})$. For the GeoVac spectral triple on round
$\sthree$, the Seeley--DeWitt coefficients of $\DCH^{2}$ are exact in
$\sqrt{\pi}\cdot\Q$ with no higher coefficients (Paper~28~\cite{paper28},
Sprint MR-B closed form, two-term-exact theorem). The spectral action
therefore lives in the M2 ring of the master Mellin engine
(Paper~18~\cite{paper18} \S III.7).

\textbf{Structural correspondence.} Both the Henningson--Skenderis
holographic Weyl anomaly and the Connes--Chamseddine spectral action
sit in the master Mellin engine's M2 sub-mechanism (Seeley--DeWitt
heat-kernel ring $\sqrt{\pi}\cdot\Q \oplus \pi^{2}\cdot\Q$). Their
correspondence is not a numerical coincidence but a structural
inheritance: both compute the same operator-algebraic object (heat-kernel
trace of $D^{2}$, evaluated via $\mathcal{M}[\mathrm{Tr}\,e^{-tD^{2}}]$
at operator order $k = 2$) from two different sides --- bulk gravity
(HS) and finite spectral triple (CC). The framework's discrete-cutoff
realization (Paper~38~\cite{paper38} propinquity convergence at rate
$4/\pi$) provides the operator-system-level realization of this
correspondence at finite cutoff.

This is the framework's most direct contact point with established
holographic computations in the M2 sector and is a corollary of the
master Mellin engine case-exhaustion theorem (Paper~32 \S VIII), not
a new computation. The substantive content is the explicit
identification of HS's bulk computation as an M2-sector observable,
which the framework's spectral action machinery on $\sthree$
reproduces structurally at finite cutoff.

\paragraph{Relation to the TT-bar cut-off AdS picture.}
Hartman, Kruthoff, Shaghoulian, and Tajdini~\cite{hartman_kruthoff_shaghoulian_tajdini2019}
showed that sphere partition functions of TT-bar-deformed large-$N$
holographic CFT$_d$ in $d \in \{2, 3, 4, 5, 6\}$ non-perturbatively
match bulk computations of AdS$_{d+1}$ at finite radial cut-off. The
parallel to the present paper's discrete-cutoff framework
($\Tcal_{\nmax}$ at finite $\nmax$ matching continuum CFT) is suggestive,
but the two cut-off structures are \emph{categorically distinct}:

\begin{itemize}
\item \textbf{TT-bar (Hartman et al.~2019):} a \emph{deformation of the
action}. The TT-bar operator $\mathcal{O}_{T\bar T}(\mu)$ modifies the
field theory's UV at finite $\mu$; the deformed CFT has a different
energy spectrum and a different action. The bulk dual is AdS at finite
radial cutoff $r_c$ related to $\mu$ via a (dimension-dependent)
explicit formula.

\item \textbf{GeoVac $\nmax$ (this paper):} a \emph{truncation of the
spectrum}, not the action. The framework's discrete spectral triple at
$\nmax$ keeps the underlying free CFT and its action intact, but
restricts the mode sum to $n \le \nmax$. As $\nmax \to \infty$, the
spectral data converges to the round-sphere CFT bit-exactly
(Theorems~\ref{thm:scalar_bitexact},~\ref{thm:dirac_bitexact}); the
intermediate $\nmax$ values are not a TT-bar-deformed CFT but
genuine spectral truncations of the original CFT.
\end{itemize}

\textbf{At the partition-function-value level,} the two regularizations
might still be related: each produces a finite number $Z_{\text{cutoff}}$
that converges to the continuum CFT partition function as the cutoff is
removed. Whether there exists a $\mu(\nmax, R)$ such that
$Z_{\nmax}^{\text{GeoVac}}(R) = Z_{\text{TT-bar}}(\mu(\nmax, R), R)$ is
an open question. A positive answer would identify the framework's
spectral cutoff with a specific TT-bar-deformation strength; a negative
answer would crystallize that the two cutoffs are structurally
incompatible at the partition-function level.

\paragraph{Structural distinction from holographic correspondences (dS/CFT and AdS/CFT).}
\label{para:category_iii_position}
The bit-exact reproduction of CFT$_{3}$-on-$\sthree$ partition function
data (Sec.~\ref{sec:partition}) places the framework in structural
adjacency with both AdS/CFT and dS/CFT, but in neither category.
Specifically, the framework's $\sthree$ substrate occupies the same
geometric position that Euclidean dS/CFT~\cite{maldacena2002} assigns to
its boundary CFT$_{3}$: in Maldacena's Euclidean sharpening of
Strominger's dS/CFT proposal~\cite{strominger2001}, the Hartle--Hawking
wavefunctional of dS$_{4}$ at $I^{+}$ equals the CFT$_{3}$-on-$\sthree$
partition function $\Psi_{\mathrm{dS}_{4}}[\phi_{0}] = Z_{\mathrm{CFT}_{3}}[\phi_{0}]$.
The KPS values reproduced bit-exactly here
(Theorems~\ref{thm:scalar_bitexact},~\ref{thm:dirac_bitexact}) are
precisely the boundary CFT data that dS/CFT would predict via this
identification.

However, the framework does \emph{not} contain a bulk dS$_{4}$ (or its
Euclidean partner $S^{4}$). The $\sthree$ is the Fock-projection image
of the Coulomb spectrum (Paper~7~\cite{paper7}), not a future-infinity
boundary of any 4D bulk; there is no Hartle--Hawking wavefunctional
construction (no functional integral over 4D geometries); and there is
no holographic dictionary mapping bulk fields to boundary operators.
The CFT$_{3}$ partition function values emerge directly from spectral
asymptotics of the discrete spectral triple, not from holographic
reconstruction. This places the framework in a third structural
category, distinct from the holographic correspondences:

\begin{center}
\begin{tabular}{lll}
\toprule
Theory & Bulk & Correspondence mechanism \\
\midrule
AdS$_{5}$/CFT$_{4}$ & AdS$_{5}$ IIB SUGRA & Holographic dictionary \\
AdS$_{4}$/CFT$_{3}$ (ABJM) & AdS$_{4}$ M-theory & Holographic dictionary \\
dS$_{4}$/CFT$_{3}$ & dS$_{4}$ & Hartle--Hawking wavefunctional \\
Chern--Simons / WZW & 3D topological & State-sum / topological \\
Connes--vS (2021) $S^{1}$ & Discrete spectral triple & Propinquity convergence \\
\textbf{GeoVac} (this paper) & \textbf{Discrete spectral triple on $\sthree$} & \textbf{Latr\'emoli\`ere propinquity convergence at rate $4/\pi$} \\
\bottomrule
\end{tabular}
\end{center}

The framework's analog of the holographic correspondence is the
Paper~38~\cite{paper38} propinquity convergence theorem
$\Lambda(\Tcal_{\nmax}, \Tcal_{\infty}) \le C_{3} \gamma_{\nmax} \to 0$.
This is a correspondence between two operator-algebraic objects (the
discrete spectral triple at finite cutoff and the continuum spectral
triple on round $\sthree$), not between two physical geometries
(bulk gravity and boundary CFT). The third category is structurally
characterized by:
\begin{itemize}
\item the ``bulk'' replaced by a discrete spectral-triple substrate
(operator-system level, finite cutoff);
\item the ``correspondence'' replaced by propinquity convergence in
the continuum limit (Paper~38's five-lemma proof);
\item boundary CFT data emerging from spectral asymptotics
(zeta-function derivatives at zero), not from holographic reconstruction.
\end{itemize}

In one sentence: \emph{the framework is the discrete spectral-triple
realization of the boundary CFT$_{3}$-on-$\sthree$ partition function
data and modular structure, distinguished from AdS/CFT and dS/CFT by
absence of any bulk gravitational theory and presence of a
propinquity-convergence ``correspondence'' between operator-algebraic
discretization levels}. This is not a new holographic correspondence
(no bulk = no holography), nor a refutation of dS/CFT (the boundary
data is consistent with what dS/CFT would predict); it is a structurally
distinct realization that produces the same boundary CFT data via a
different mechanism. The bulk-side identifications named in
Prop.~\ref{prop:bulk_blocked} for AdS$_{4}$ remain unreachable from the
framework's existing infrastructure for the same reasons (no bulk
geometric machinery in \texttt{geovac/}). The framework's contribution
is to provide a non-holographic operationalization of the
boundary-CFT-on-$\sthree$ data, structurally distinct from both AdS/CFT
and dS/CFT.

\paragraph{Areas yet to be explored.}
Beyond the immediate catalogue targets, several adjacent directions
extend the framework's reach into established physics:

\begin{itemize}
\item \textbf{Holographic central charges} ($a$- and $c$-anomalies in
even $d$). Henningson-Skenderis~\cite{henningson_skenderis1998} and
follow-ups give closed forms for the holographic dual of free CFT
central charges. The framework's spectral-action machinery
(Paper~28~\cite{paper28}) potentially provides a discrete-cutoff
reproduction; the connection to Paper~24's Bargmann-Segal $S^{5}$ is
unexplored.

\item \textbf{Entanglement-entropy universal terms.} The continuum
Casini-Huerta-Myers~\cite{casini_huerta_myers2011} entanglement
entropy of a ball region has a universal subleading constant equal
to $-F$. Connecting this to the framework's wedge KMS structure
(Sec.~\ref{sec:wedge_kms}) and to the master Mellin engine
decomposition would strengthen the F-as-engine-signature reading.

\item \textbf{Higher-rank Lie groups beyond $\mathrm{SU}(2)$.} The
Paper~40 unified GH-convergence theorem~\cite{paper40} extends the
propinquity machinery to all simple compact Lie groups. A natural
catalogue extension is CFT-on-$G$ partition functions for various
$G$ (heat-kernel zeta on $\mathrm{SU}(3)$, $\mathrm{Sp}(2)$, etc.)
where the framework's GH-rate constant $4/\pi$ is universal but the
specific transcendental ring may differ.

\item \textbf{Spectral-action coefficients.} The
Connes-Chamseddine~\cite{chamseddine_connes1997} spectral action
$\mathrm{Tr}\,f(D/\Lambda)$ has heat-kernel expansion coefficients
that classify the gravitational + matter content of an
almost-commutative spectral triple. These coefficients are exactly
the master Mellin engine M2 sector. The catalogue could include
M-engine readings of standard spectral-action calculations on
$\sthree$, $S^{5}$, and higher.

\item \textbf{Lorentzian-side observables.} The catalogue here is
boundary-side Riemannian. The Lorentzian extension (Papers~42-49)
provides modular Hamiltonian, Pythagorean orthogonality, and
twin-paradox-as-quantum-information structures. Cataloguing
Lorentzian-side CFT observables (thermal partition functions on
$\sthree \times \mathbb{R}_{t}$, Rindler-wedge entropies, etc.) is a
parallel direction.
\end{itemize}

Each catalogue entry either (i) verifies the master Mellin engine on
independent physics data, (ii) extends a structural cancellation
theorem (like log~3 cancellation on $S^{5}$,
Theorem~\ref{thm:log3_cancellation}), or (iii) identifies a wall
where the simple ring fails (like the $S^{7}$ scalar negative,
Remark~\ref{rem:general_d_conjecture}). Each outcome is paper-worthy
in its own right; the cumulative catalogue would constitute a
structural classification of free-field CFT-on-sphere partition
functions via the GeoVac spectral triple.

\subsection{Place in the GeoVac math.OA series}

This paper is the twelfth math.OA standalone in the GeoVac series,
following Papers~38, 39, 40, 42, 43, 44, 45, 46, 47, 48, 49.
The previous eleven established operator-algebraic legitimacy of the
framework's discrete-spectral-triple truncation, the propinquity
convergence theory at various levels (Riemannian, tensor-product,
universal compact-Lie-group, K$^{+}$-weak-form Lorentzian,
strong-form Lorentzian, norm-resolvent non-compact, K$^{+}$-weak-form
synthetic Lorentzian via Mondino--S\"amann bridge, strong-form
synthetic Lorentzian via OSLPLS). The present paper is the
\emph{first connecting that math.OA arc to a physics literature}
--- specifically the CFT-on-sphere partition function literature
of~\cite{cardy1988, jafferis_klebanov_pufu_safdi2011,
klebanov_pufu_safdi2011, beccaria_tseytlin2017,
lei_van_leuven2024} --- with a load-bearing structural finding
(the master Mellin engine verification on independent physics
data) plus three bit-exact / formal correspondences (partition
function, wedge boundary dimension, CHM modular Hamiltonian).

\section*{Acknowledgments}

This work was developed in the GeoVac project's AI-augmented
agentic workflow with Anthropic's Claude.

\begin{thebibliography}{99}

\bibitem{paper7} J.~Loutey,
\emph{Paper~7: Dimensionless Vacuum},
GeoVac internal preprint, 2026.

\bibitem{paper18} J.~Loutey,
\emph{Paper~18: Exchange Constants and Transcendental Taxonomy},
GeoVac internal preprint, 2026.

\bibitem{paper22} J.~Loutey,
\emph{Paper~22: Angular Sparsity Theorem on $\sthree$},
GeoVac internal preprint, 2026.

\bibitem{paper23} J.~Loutey,
\emph{Paper~23: Nuclear Shell Model on the Bargmann--Segal Lattice},
GeoVac internal preprint, 2026.

\bibitem{paper29} J.~Loutey,
\emph{Paper~29: Ramanujan Hopf Graphs and Wick-Rotation Scope},
GeoVac internal preprint, 2026.

\bibitem{paper32} J.~Loutey,
\emph{Paper~32: The GeoVac Spectral Triple},
GeoVac internal preprint, 2026.

\bibitem{paper35} J.~Loutey,
\emph{Paper~35: Time as Projection},
GeoVac internal preprint, 2026.

\bibitem{paper38} J.~Loutey,
\emph{Paper~38: SU(2) Propinquity Convergence on the Camporesi--Higuchi Spectral Triple},
GeoVac internal preprint, 2026.

\bibitem{paper28} J.~Loutey,
\emph{Paper~28: QED on $S^{3}$},
GeoVac internal preprint, 2026.

\bibitem{paper40} J.~Loutey,
\emph{Paper~40: Unified Propinquity Convergence on Compact Lie Groups},
GeoVac internal preprint, 2026.

\bibitem{paper42} J.~Loutey,
\emph{Paper~42: Tomita--Takesaki Modular Structure and Four-Witness Wick Rotation},
GeoVac internal preprint, 2026.

\bibitem{paper43} J.~Loutey,
\emph{Paper~43: Lorentzian Extension of the Four-Witness Theorem},
GeoVac internal preprint, 2026.

\bibitem{paper44} J.~Loutey,
\emph{Paper~44: Operator-System Substrate for Lorentzian Propinquity},
GeoVac internal preprint, 2026.

\bibitem{paper45} J.~Loutey,
\emph{Paper~45: K$^{+}$-Weak-Form Lorentzian Propinquity Convergence},
GeoVac internal preprint, 2026.

\bibitem{paper46} J.~Loutey,
\emph{Paper~46: Strong-Form Lorentzian Propinquity Convergence},
GeoVac internal preprint, 2026.

\bibitem{paper47} J.~Loutey,
\emph{Paper~47: Norm-Resolvent Convergence to the Non-Compact Lorentzian Krein Spectral Triple},
GeoVac internal preprint, 2026.

\bibitem{paper48} J.~Loutey,
\emph{Paper~48: Krein--Mondino--S\"amann Bridge},
GeoVac internal preprint, 2026.

\bibitem{paper49} J.~Loutey,
\emph{Paper~49: OSLPLS Strong-Form Bridge and Twin Paradox as Quantum Information},
GeoVac internal preprint, 2026.

\bibitem{fock1935} V.~Fock,
\emph{Zur Theorie des Wasserstoffatoms},
Z.~Phys.\ \textbf{98} (1935), 145--154.

\bibitem{camporesi_higuchi1996} R.~Camporesi and A.~Higuchi,
\emph{On the eigenfunctions of the Dirac operator on spheres and real hyperbolic spaces},
J.~Geom.~Phys.\ \textbf{20} (1996), 1--18.

\bibitem{connes_vs2021} A.~Connes and W.~D.~van Suijlekom,
\emph{Spectral truncations in noncommutative geometry and operator systems},
Comm.~Math.~Phys.\ \textbf{383} (2021), 2021--2067, arXiv:2004.14115.

\bibitem{cardy1988} J.~L.~Cardy,
\emph{Operator content of two-dimensional conformally invariant theories},
Nucl.~Phys.~B \textbf{270} (1986), 186--204.

\bibitem{klebanov_pufu_safdi2011} I.~R.~Klebanov, S.~S.~Pufu, and B.~R.~Safdi,
\emph{F-theorem without supersymmetry},
JHEP \textbf{10} (2011) 038, arXiv:1105.4598.

\bibitem{jafferis_klebanov_pufu_safdi2011} D.~L.~Jafferis, I.~R.~Klebanov, S.~S.~Pufu, and B.~R.~Safdi,
\emph{Towards the F-theorem: $N = 2$ field theories on the three-sphere},
JHEP \textbf{06} (2011) 102, arXiv:1103.1181.

\bibitem{beccaria_tseytlin2017} M.~Beccaria and A.~A.~Tseytlin,
\emph{$\mathcal{N} = 4$ conformal supergravity: the complete one-loop dilatation operator},
JHEP \textbf{04} (2017) 100, arXiv:1702.02325.

\bibitem{lei_van_leuven2024} J.~Lei and W.~D.~van Leuven,
\emph{Modularity in $d > 2$ free conformal field theory},
arXiv:2406.01567, 2024.

\bibitem{henningson_skenderis1998} M.~Henningson and K.~Skenderis,
\emph{The holographic Weyl anomaly},
JHEP \textbf{07} (1998) 023, arXiv:hep-th/9806087.

\bibitem{hartman_kruthoff_shaghoulian_tajdini2019} T.~Hartman, J.~Kruthoff,
E.~Shaghoulian, and A.~Tajdini, \emph{Holography at finite cutoff with a
$T^{2}$ deformation}, JHEP \textbf{03} (2019) 004, arXiv:1902.10893.

\bibitem{chamseddine_connes1997} A.~H.~Chamseddine and A.~Connes,
\emph{The spectral action principle},
Comm.~Math.~Phys.\ \textbf{186} (1997), 731--750.

\bibitem{casini_huerta_myers2011} H.~Casini, M.~Huerta, and R.~C.~Myers,
\emph{Towards a derivation of holographic entanglement entropy},
JHEP \textbf{05} (2011) 036, arXiv:1102.0440.

\bibitem{bisognano_wichmann1976} J.~J.~Bisognano and E.~H.~Wichmann,
\emph{On the duality condition for quantum fields},
J.~Math.~Phys.\ \textbf{17} (1976), 303--321.

\bibitem{ryu_takayanagi2006} S.~Ryu and T.~Takayanagi,
\emph{Holographic derivation of entanglement entropy from AdS/CFT},
Phys.~Rev.~Lett.\ \textbf{96} (2006), 181602, arXiv:hep-th/0603001.

\bibitem{jafferis_lewkowycz_maldacena_suh2016} D.~L.~Jafferis, A.~Lewkowycz, J.~Maldacena, and S.~J.~Suh,
\emph{Relative entropy equals bulk relative entropy},
JHEP \textbf{06} (2016) 004, arXiv:1512.06431.

\bibitem{pastawski_yoshida_harlow_preskill2015} F.~Pastawski, B.~Yoshida, D.~Harlow, and J.~Preskill,
\emph{Holographic quantum error-correcting codes: toy models for the bulk/boundary correspondence},
JHEP \textbf{06} (2015) 149, arXiv:1503.06237.

\bibitem{debug:bh_phase0} GeoVac internal,
\emph{Sprint BH-Phase0 diagnostic memo},
\texttt{debug/bh\_phase0\_diagnostic\_memo.md}, 2026-05-22.

\bibitem{debug:sprint_td_track5} GeoVac internal,
\emph{Sprint TD Track 5 PSLQ-disjoint correlation entropy},
\texttt{debug/sprint\_td\_track5\_memo.md}, 2026-05-09.

\bibitem{debug:external_input_three_class_partition} GeoVac internal,
\emph{External-input three-class partition},
\texttt{memory/external\_input\_three\_class\_partition.md}, 2026-05-25.

\bibitem{strominger2001} A.~Strominger,
\emph{The dS/CFT correspondence},
JHEP \textbf{10} (2001) 034, arXiv:hep-th/0106113.

\bibitem{maldacena2002} J.~Maldacena,
\emph{Non-Gaussian features of primordial fluctuations in single field inflationary models},
JHEP \textbf{05} (2003) 013, arXiv:astro-ph/0210603.

\end{thebibliography}

\end{document}
